\chapter{光的干涉~几何光学}
\setcounter{example_counter}{1}

\section{光的干涉}

\subsection{光程}

\subsubsection{相干光}

\subsubsection{光在介质中的传播}

光学介质的折射率$n$，是光在真空中的传播速度与光在介质中传播速度之比。

\begin{formula}
    $n=\dfrac{c}{v}$
\end{formula}

\note{光在真空中传播速度最快，所以折射率不会小于1.}

两种介质相比，折射率$n$较大的称为光密介质，折射率较小的称为光疏介质。

\subsubsection{光程与光程差}

折射率$n$与光在介质中的路程$r$的乘积称为光程。

\begin{formula}
    光程$=nr$
\end{formula}

光在折射率$n$的介质中传播的速率是$v=\dfrac{c}{n}$，因此在时间$t$内，光在介质中传播的路程$=nr=ct$.~这说明在相同的时间内，光在不同介质内传播的几何路程不同，但光程是相同的。

\subsubsection{光程与相位差}

光程是按相位变化折算到真空中的路程，利用光程差可以直接得到相位差。

\begin{formula}
    相位差$=\dfrac{2\pi}{\lambda}\times$光程差
\end{formula}

但要注意，光从光疏介质射入光密介质，所产生的反射波还会出现半波损失，此时相位差还要$\pm \pi$

\begin{example}{光程差与相位差}
    $S_1,S_2$是两个相干光源，相位相同，到P点距离分别为$r_1,r_2$，路径$S_1P$垂直穿过一块厚度为$t_1$，折射率为$n_1$的介质板，路径$S_2P$垂直穿过一块厚度为$t_2$，折射率为$n_2$的介质板，其余部分可看做真空。(1)~求这两条路径的光程差；(2)~若光在真空中的波长为$\lambda$，求两束光在P点相遇时的相位差。（310题）
    \tcblower
    如图所示，平行单色光垂直照射到薄膜上，薄膜的厚度为$e$，薄膜及周围介质的折射率如图所示。在下列两种情况下，求经上下两表面反射的光在相遇时的相位差。(1)~$n_3=1.0$，(2)~$n_3=1.5$ （第311题）
\end{example}

\subsection{双缝干涉}

\subsection{薄膜干涉}

\subsubsection{等厚条纹}

\subsubsection{等倾条纹}


\section{几何光学}

\subsection{反射与折射}

\subsection{球面镜}

\subsection{薄透镜}